% !TITLE 玉楼春·晚妆初了明肌雪
% !AUTHOR 李煜
% !DYNASTY 五代十国
% !GENRE 词
% !ORDER 2
% !EDITION 《南唐二主词笺注》，王仲闻校订、陈书良与刘娟笺注，中华书局（2013年版为候选）
% !EDITION_STATUS pending

\begin{poem}
晚妆\textsuperscript{1}初了明肌雪\textsuperscript{2}，春殿嫔娥\textsuperscript{3}鱼贯列\textsuperscript{4}。
凤箫\textsuperscript{5}吹断水云闲，重按霓裳\textsuperscript{6}歌遍彻\textsuperscript{7}。

临风谁更飘香屑\textsuperscript{8}，醉拍阑干情味切\textsuperscript{9}。
归时休放烛花红\textsuperscript{10}，待踏马蹄清夜月\textsuperscript{11}。
\end{poem}

\begin{annotations}
\notetext{1}{晚妆：晚上的妆扮。不是早起梳妆，而是晚上宴会前重新梳妆——可见宫廷生活的奢华。}
\notetext{2}{明肌雪：肌肤明亮如雪。晚妆之后，宫女们容光焕发，这是亡国前的繁华景象。}
\notetext{3}{嫔娥：指宫女和妃嫔。嫔（pín），帝王妃子的一个等级；娥，美女的泛称。}
\notetext{4}{鱼贯列：像鱼一样一个接一个排列整齐。宫女们依次走入殿中，阵容整齐壮观。}
\notetext{5}{凤箫：排箫（一种管乐器），因声音清亮如凤鸣而得名。也有人解释为做成凤凰形状的箫。}
\notetext{6}{霓裳：指《霓裳羽衣曲》，唐玄宗时代最著名的大型乐舞。南唐宫中反复演奏此曲，可见李煜对盛唐文化的追慕。}
\notetext{7}{遍彻：从头到尾完整演奏。“遍”是大曲的一个段落，整部大曲从头到尾演完。}
\notetext{8}{香屑：碾碎的香料粉末。有人在风中抛洒香粉，空气中弥漫着芬芳——宴会不止有音乐，还有嗅觉的享受。}
\notetext{9}{情味切：情感深切动人。喝醉了拍着栏杆，陷入最深的情致中。}
\notetext{10}{烛花红：红烛的光。回去时不要点亮烛火——因为月光太美了，不想被灯光破坏。}
\notetext{11}{清夜月：清朗的夜月。踏着马蹄声，沐浴在月光下归去——通宵欢宴后，最后一幕是策马踏月。这是李煜笔下最美的场景之一，也是他后半生再也回不去的记忆。}
\end{annotations}

\begin{appreciation}
这是李煜夜宴的又一夜。写宴饮，他不再铺排酒馔，只拣声色两样来写，乐是《霓裳》，色是月色，最终压在词尾的，竟是一地清光。

“晚妆初了明肌雪”，是晚妆，不是晨妆。晨起梳洗过了，赴宴前还要重新理一遍妆容，这一层讲究，就把宫廷的气派点足了。妆成之后肌肤明净如雪，宫女们“鱼贯列”，一个挨一个走进殿来，像游鱼成行，连队形也是一种排场。

“凤箫吹断水云闲，重按霓裳歌遍彻。”凤箫是排箫，声音清亮如凤鸣；吹断了水与云之间的闲气，是说乐声悠长，把天边的云和水都吹得安闲下来。“重按霓裳”是按谱重新演奏《霓裳羽衣曲》，“遍彻”则是整部大曲从第一叠一直奏到入破，一遍不落。

《霓裳羽衣曲》本是唐玄宗时代最盛大的法曲，安史之乱后散佚，只剩残谱。据陆游《南唐书》所载，南唐宫中得到了这份残谱，由精于琵琶的昭惠国后按谱整理去讹补缺，才重新奏出了开元、天宝的遗音。这段记载究竟有几分可信今人已难断言，但有一点是清楚的，李煜对盛唐的追慕是从一曲旧乐里一点一点找回来的。盛唐的气象早已散尽，曲子却还在南唐的深宫一遍遍响起。

“临风谁更飘香屑，醉拍阑干情味切。”有人在风里抛撒碾碎的香屑，让香气随风散开，连鼻子也一并伺候到了；酒意上来，拍着栏杆，动了真情。

最妙的是收尾两句：“归时休放烛花红，待踏马蹄清夜月。”回宫的路上不必点灯，就让马蹄踏着一地月光走。热闹了一整夜，他最后舍不得的偏是那一点冷清，众人皆醉，他却惦记着月色。这是李煜早期词里最好的一笔。

《鹿鸣》里的宴饮，鼓瑟吹笙，宾主尽欢，热闹里全是一团和气；李煜的夜宴排场大了许多，主角却只剩他一个人的兴致。同样是宴，一个写的是人与人之间的和乐，一个写的是人与声色的独对。

“待踏马蹄清夜月”这个画面，日后他再也没能走回去。汴京的月亮照样照着，只是马蹄声没有了。而他要月光、不要烛火的那点偏爱，其实早已露出了征兆：这个人的心里原就藏着一条只属于他自己的清冷的路。
\end{appreciation}
